\section*{Standard-Integrale und Volumenelemente}

\subsection*{1. Kugelkoordinaten (3D)}
\textbf{Allgemeines Volumenelement:}
\[
dV = r^2 \sin(\theta) \, d\theta \, d\phi \, dr
\]
\textbf{Kugelvolumen:}
\[
V = \int_{0}^{R} \int_{0}^{\pi} \int_{0}^{2\pi} r^2 \sin(\theta) \, d\phi \, d\theta \, dr = \frac{4}{3}\pi R^3
\]
\textbf{Shortcut (Kugelschale bei reiner $r$-Abhängigkeit):}
\begin{formulabox}
	\[
	dV = 4\pi r^2 \, dr
	\]
\end{formulabox}

\subsection*{2. Zylinderkoordinaten (3D)}
\textbf{Allgemeines Volumenelement:}
\[
dV = r \, dr \, d\phi \, dz
\]
\textbf{Zylindervolumen:}
\[
V = \int_{0}^{h} \int_{0}^{2\pi} \int_{0}^{R} r \, dr \, d\phi \, dz = \pi R^2 h
\]
\textbf{Shortcut (Zylinderschale bei reiner $r$-Abhängigkeit):}
\begin{formulabox}
	\[
	dV = h \cdot 2\pi r \, dr
	\]
\end{formulabox}

\subsection*{3. Ebene Polarkoordinaten (2D)}
\textbf{Allgemeines Flächenelement:}
\[
dA = r \, dr \, d\phi
\]
\textbf{Shortcut (Ring bei reiner $r$-Abhängigkeit):}
\begin{formulabox}
	\[
	dA = 2\pi r \, dr
	\]
\end{formulabox}

\subsection*{4. Kartesische Koordinaten}
\textbf{Volumen (3D):}
\[
dV = dx \, dy \, dz
\]
\textbf{Fläche (2D):}
\[
dA = dx \, dy
\]

\subsection*{5. Wichtige Linienelemente ($ds$) und Oberflächen ($A$)}
\begin{itemize}
	\item \textbf{Kreisbogen:}
	\begin{formulabox}
		\[
		ds = r \, d\phi
		\]
	\end{formulabox}
	
	\item \textbf{Kugeloberfläche:}
	\[
	A = 4\pi r^2
	\]
	
	\item \textbf{Kugelzone (fester Radius $R$):}
	\[
	dA = 2\pi R^2 \sin(\theta) \, d\theta
	\]
	
	\item \textbf{Zylindermantel:}
	\[
	A = 2\pi r h
	\]
\end{itemize}

\subsection*{6. Vektorelemente für Felder (z.\,B. Maxwell-Gleichungen)}
\textbf{Vektor-Flächenelement:}
\[
d\vec{A} = dA \cdot \hat{n}
\]
(Der Normalenvektor $\hat{n}$ zeigt bei geschlossenen Hüllen immer nach aussen!) \\[1ex]
\textbf{Vektor-Linienelement:} $d\vec{l}$ \\
Zeigt in die gewählte Integrationsrichtung oder (bei Leitern) in Richtung der Stromdichte $\vec{j}$.